\documentclass[12pt,a4paper]{report}

% ── Encoding & Fonts ──
\usepackage[utf8]{inputenc}
\usepackage[T1]{fontenc}
\usepackage{lmodern}

% ── Page Layout ──
\usepackage[margin=1in,headheight=15pt]{geometry}
\usepackage{setspace}
\onehalfspacing

% ── Language ──
\usepackage[english]{babel}

% ── Graphics & Colors ──
\usepackage{graphicx}
\usepackage[dvipsnames,svgnames,table]{xcolor}

% ── Math ──
\usepackage{amsmath,amssymb,amsfonts}

% ── Tables ──
\usepackage{booktabs}
\usepackage{tabularx}
\usepackage{longtable}
\usepackage{multirow}

% ── Lists ──
\usepackage{enumitem}

% ── Code Listings ──
\usepackage{listings}
\lstset{
  basicstyle=\ttfamily\small,
  keywordstyle=\color{blue}\bfseries,
  commentstyle=\color{gray},
  stringstyle=\color{red},
  breaklines=true,
  frame=single,
  backgroundcolor=\color{gray!5},
  numbers=left,
  numberstyle=\tiny\color{gray},
  numbersep=8pt,
  tabsize=4,
  showstringspaces=false,
}

% ── Highlighting ──
\usepackage{ulem}
\normalem
% Custom highlight command (soul.sty not available)
\newcommand{\hl}[1]{\colorbox{yellow}{#1}}

% ── Hyperlinks ──
\usepackage[colorlinks=true,linkcolor=NavyBlue,citecolor=Green,urlcolor=RoyalBlue]{hyperref}

% ── Headers & Footers ──
\usepackage{fancyhdr}
\pagestyle{fancy}
\fancyhf{}
\fancyhead[L]{\leftmark}
\fancyhead[R]{\thepage}
\renewcommand{\headrulewidth}{0.4pt}

% ── Chapter Style ──
\usepackage{titlesec}
\titleformat{\chapter}[display]
  {\normalfont\huge\bfseries}
  {\chaptertitlename\ \thechapter}{20pt}{\Huge}
\titlespacing*{\chapter}{0pt}{-20pt}{30pt}

% ── Caption Styling ──
\usepackage{caption}
\captionsetup{font=small,labelfont=bf}

% ── Bibliography ──
\usepackage[numbers,sort&compress]{natbib}

% ── Misc ──
\usepackage{parskip}
\usepackage{float}

% ════════════════════════════════════════════════
%  DOCUMENT METADATA — Update these as needed
% ════════════════════════════════════════════════
\title{\textbf{CFD Setup Instructions}}
\author{Yash}
\date{\today}

\begin{document}

% ── Title Page ──
\maketitle
\thispagestyle{empty}

% ── Table of Contents ──
\tableofcontents
\newpage

% ════════════════════════════════════════════════
%  CHAPTERS — Content will be added below
% ════════════════════════════════════════════════

% ────────────────────────────────────────────────
%  CHAPTER 1
% ────────────────────────────────────────────────
\chapter{Geometry Creation in DesignModeler}

% ── Section 1 ──
\section{Prepare DesignModeler}

\textbf{From Workbench}

Right-click the Geometry cell $\rightarrow$ New Geometry $\rightarrow$ DesignModeler.

When prompted, select the desired length unit (\hl{\textbf{Millimeter is recommended}}).

\subsection*{Change to ``Inches'' temporarily (optional)}

The main pipe is specified as \textbf{18 inch} in the paper.

You can set the unit to \textbf{Inch (Units $\rightarrow$ Inch)} to enter \textbf{18 directly}, then switch back to \textbf{Millimeter} for the rest of the dimensions.

If you prefer to stay in mm, simply enter the metric equivalents.


% ── Section 2 ──
\section{Create the Main (Horizontal) Pipe}

\subsection*{Select the sketching plane}

In the Tree Outline, expand \textbf{XYPlane $\rightarrow$ click Sketching}.

\subsection*{Draw the main pipe cross-section}

In the Draw toolbox, choose \textbf{Circle}.

Click the origin (0,0) to place the centre, then drag out a circle.

In the Details window under Dimensions set the \textbf{Diameter to 460\,mm (or 18\,in)}.

\textbf{Why?} The paper states: \textit{``In the UK, the current natural gas network consists of 18 inch (460\,mm) diameter pipes''}.

\subsection*{Extrude the main pipe}

Switch to \textbf{Modeling (top of the Tree Outline)}.

Select the circle face you just sketched.

Click \textbf{Extrude} in the Create toolbar.

In the Details window:

\begin{itemize}
  \item \textbf{Extent Type:} Fixed
  \item \textbf{Depth (>0):} 10\,m ($\approx$ 10{,}000\,mm) for the upstream section
  \item \textbf{Direction:} Reversed if you want the origin at the inlet face (optional)
\end{itemize}

Click \textbf{Generate}.

\textbf{Why 10\,m?} The upstream length should be at least \textbf{10 pipe diameters ($10 \times 460$\,mm)} to allow the main flow to become fully developed before the branch. The paper does not specify a length, but this is a standard CFD practice.

\subsection*{Rename the body (optional)}

In the Tree Outline right-click \textbf{Extrude\,1 $\rightarrow$ Rename $\rightarrow$ MainPipe}.


% ── Section 3 ──
\section{Create the Branch (Vertical) Pipe}

\subsection*{Define a new sketch plane at the branch location}

The branch should be centred at a distance \textbf{5\,m from the main inlet} (i.e., half of the 10\,m upstream length).

In the Tree Outline click \textbf{New Plane $\rightarrow$ Plane from Face}.

Select the outer cylindrical face of the main pipe and \textbf{offset the plane by 5\,m along the main pipe axis (Z-direction)}.

Alternatively, use the \textbf{XYPlane} and extrude the branch from there, then move it into position.

\subsection*{Sketch the branch cross-section}

On the newly created plane, enter \textbf{Sketching}.

Draw a circle centred on the main pipe axis.

Set its \textbf{Diameter to 115\,mm} (since $d_2/d_1 = 1/4$ and $d_1 = 460$\,mm).

\textbf{Why?} The paper examines \textit{``geometrical differences of $d_2/d_1 = 1/2,\; 1/4$''}. For the $1/4$ case, the branch diameter is \textbf{$460/4 = 115$\,mm}.

\subsection*{Extrude the branch}

Go back to \textbf{Modeling}, select the branch circle.

Click \textbf{Extrude}.

\begin{itemize}
  \item \textbf{Extent Type:} Fixed
  \item \textbf{Depth (>0):} 1\,m ($10 \times$ branch diameter, a typical length to ensure the branch flow develops before entering the main pipe)
  \item \textbf{Direction:} Normal (upwards for top-side injection)
\end{itemize}

Click \textbf{Generate}.

Rename to \textbf{BranchPipe}.

\subsection*{Position the branch (if not already centred at the correct Z-coordinate)}

Use \textbf{Body Transform $\rightarrow$ Translate} to move the branch so that its centre lies exactly at the desired distance from the main inlet.


% ── Section 4 ──
\section{Merge the Two Pipes}

\subsection*{Select both bodies}

In the Tree Outline, select \textbf{MainPipe and BranchPipe (Ctrl+click)}.

\subsection*{Unite them}

Click \textbf{Create $\rightarrow$ Boolean $\rightarrow$ Unite}.

The two bodies become a \textbf{single fluid volume}.

\textbf{Why?} You need a \textbf{watertight, continuous domain for CFD}. The Unite operation removes the internal overlapping face, leaving only the outer walls.

\subsection*{Check that the geometry is a single body}

In the Tree Outline you should see \textbf{one Part containing one Body}.


% ── Section 5 ──
\section{Named Selections (for later use in Fluent)}

Named selections make it easy to apply boundary conditions without manually picking faces each time.

\subsection*{Select the main inlet face}

Pick the circular face at the upstream end of the main pipe.

Right-click $\rightarrow$ Named Selection $\rightarrow$ enter \textbf{main\_inlet}.

\subsection*{Select the branch inlet face}

Pick the circular face at the bottom of the branch.

Right-click $\rightarrow$ Named Selection $\rightarrow$ enter \textbf{branch\_inlet}.

\subsection*{Select the outlet face}

Pick the circular face at the downstream end of the main pipe.

Right-click $\rightarrow$ Named Selection $\rightarrow$ enter \textbf{outlet}.

\subsection*{Select all wall faces}

Use \textbf{Ctrl+click} to select every remaining face (all cylindrical surfaces, the T-junction inner corner, etc.).

Right-click $\rightarrow$ Named Selection $\rightarrow$ enter \textbf{wall}.

\textbf{Tip:} The wall selection can be later split into separate groups (e.g., \textbf{upper\_wall, lower\_wall}) in Fluent if you need more detailed post-processing.


% ── Section 6 ──
\section{Final Checks}

\begin{itemize}
  \item \textbf{Dimensions:} Use the Measure tool to verify the main pipe diameter (\textbf{460\,mm}), branch diameter (\textbf{115\,mm}), and lengths
  \item \textbf{Watertightness:} Look at the Tree Outline -- if the Unite operation was successful you will see only one body. No gaps should be visible when you rotate the view
  \item \textbf{Named selections:} Right-click any selection and choose \textbf{Show} to ensure the correct faces are highlighted
\end{itemize}


% ── Section 7 ──
\section{Exit and Save}

File $\rightarrow$ \textbf{Save As\ldots} $\rightarrow$ give the geometry a meaningful name\\(e.g., \textbf{T\_Junction\_Hydrogen\_Base.wbpj}).

Close DesignModeler (\textbf{the geometry is automatically transferred to the Geometry cell in Workbench}).


% ────────────────────────────────────────────────
%  CHAPTER 2
% ────────────────────────────────────────────────
\chapter{Literature Review}

\section{The Foundation and The Problem}

\subsection{The Big Picture Goal}

To fight climate change, there is a major drive to reduce carbon dioxide emissions, especially from domestic heating.

The proposed solution is to create a \textit{``blend''} by injecting \textbf{green hydrogen} into existing natural gas pipelines, targeting a hydrogen concentration of \textbf{5\% to 20\% by volume}.

At these percentages, consumers wouldn't need to change their home boilers.

\subsection{The Core Engineering Challenge}

Currently, the UK gas network is primarily made of steel and only allows a maximum hydrogen concentration of \textbf{0.1\%}.

This limit must be safely raised to deliver the new hydrogen blends.

However, blending hydrogen into natural gas creates two directly conflicting problems:

\begin{description}[style=nextline, font=\normalfont\bfseries, leftmargin=1.5em]

  \item[The Energy Density Problem:]
  Pure hydrogen holds a lot of energy per kilogram, but its density is \textbf{9 times lower} than natural gas. This means a hydrogen-blended gas contains less energy per unit of volume. To deliver the same amount of heating energy to a home, you must either pump the gas faster or increase the pipeline pressure.

  \item[The Material Safety Problem:]
  Hydrogen causes carbon steel to become brittle---a process known as \textbf{embrittlement}. This reduces the steel's ``toughness,'' making it much easier for existing microscopic defects or welds to turn into critical, fast-growing cracks under stress. Because of this danger, safety codes typically require that you \textit{decrease} the pipeline pressure.

  \item[The Conflict:]
  Energy demand requires \textbf{higher} pressures, but safety regulations require \textbf{lower} pressures.

\end{description}

\subsection{The Mixing Problem (The Focus of the Paper)}

Hydrogen is typically injected into the main natural gas line through a \textbf{``T-junction''} (a pipe connecting at a 90-degree angle).

Because hydrogen is so light compared to natural gas, it lacks the momentum to penetrate deeply into the main gas flow.

Instead of mixing smoothly, the buoyant hydrogen floats to the top of the pipe. This is called \textbf{stratification}.

The authors' computer simulations show that injecting hydrogen from the top (\textbf{top-side injection}) results in poor mixing, leaving extremely high concentrations of hydrogen (\textbf{over 40\%}) clinging to the upper wall of the pipe for long distances.

This is dangerous because exposing the top of the steel pipe to such high concentrations of hydrogen rapidly accelerates the \textbf{embrittlement} process.

\medskip
\noindent\fbox{\parbox{\dimexpr\textwidth-2\fboxsep-2\fboxrule}{%
\textbf{Key Finding previewed in the abstract:} Injecting the hydrogen from the bottom (\textbf{under-side injection}) forces the light hydrogen to rise \textit{through} the natural gas, which promotes much better mixing and reduces the dangerous concentration on the pipe walls.}}


\end{document}
